% P23
\section{Introduction}\label{sec:introduction}

Spam filtering is commonly offered as a cloud service: the provider maintains a model trained on large amounts of labelled data, and clients send their messages to be classified.
The price is privacy, since the provider sees the content of every message, including the legitimate ones, which are the large majority.
Homomorphic encryption removes this trade-off for models that can be computed on encrypted data: the client sends its message encrypted, the server computes the classifier's score without decrypting it, and only the client learns the result.

Fully homomorphic schemes can evaluate arbitrary functions but are expensive.
Many classifiers, however, only need a linear function of the input, which an additively homomorphic scheme such as Paillier~\cite{paillier1999} computes directly.
Multinomial naive Bayes is one of them: for two classes, its decision is the sign of a log-odds score that is an inner product of the features with fixed weights plus a constant (\cref{sec:background}).
Bost et al.~\cite{bost2015} protect both the model and the data and need interactive comparison protocols for that; we keep the model public to the server, as in a typical cloud service, which reduces classification to a single round trip with one encrypted score.

In this work, we build and evaluate such a privacy-preserving spam filter for SMS messages.
Our contributions are:
\begin{itemize}
  \item a protocol in which the server computes the naive Bayes log-odds on Paillier-encrypted TF-IDF features and returns one ciphertext per message, with a fixed-point encoding that makes the encrypted computation exact and hides which features are non-zero (\cref{sec:system-design});
  \item an evaluation on the SMS Spam Collection~\cite{smsspam} with duplicates removed, showing that the encrypted classifier reaches the same quality as its plaintext counterpart, \SI{98.07}{\percent} accuracy, with identical predictions and bit-exact scores (\cref{sec:setup,sec:results});
  \item a cost analysis for key lengths from \num{1024} to \num{3072} bits and $500$ to $1000$ features, which together with the accuracy for $100$ to $1000$ features motivates the number of features (\cref{sec:results});
  \item an empirical security analysis against an honest-but-curious server that can encrypt chosen messages, which shows that a classifier on ciphertexts does not beat the majority-class baseline, while encrypting only the non-zero features would leak almost as much as the plaintext, and reusing the encryption randomness or the floating-point encoding of \texttt{phe} would let the server read features directly (\cref{sec:threat-model,sec:security}).
\end{itemize}

The implementation uses the \texttt{phe} library~\cite{phe} and scikit-learn, and all results can be reproduced with a single command.
